\documentclass[10pt,a4paper]{amsart}
\usepackage[margin=19mm]{geometry}
\usepackage{amsmath,amssymb,mathtools}
\usepackage[scr=rsfs,cal=euler]{mathalfa}
\usepackage{xfrac}
\usepackage[unicode,hidelinks,bookmarks=false]{hyperref}
\newcommand{\ie}{i.e.\ }
\newcommand{\cf}{cf.\ }
\newcommand{\resp}{respectively\ }
\newcommand{\Subsection}{\subsection}
\providecommand{\nobreakdash}{\nobreak\hbox{-}}
\setlength{\emergencystretch}{2em}
\multlinegap0pt
\usepackage{amssymb}
\newtheorem{lemma}{Lemma}
\newcommand{\E}{\mathbb E}
\newcommand{\Hh}{\mathcal H}
\newcommand{\X}{\mathcal X}
\title{A simple geometric proof for\\ the characterisation of e-merging functions}
\author{Eugenio Clerico}
\date{Source: arXiv:2512.09708v3 (CC BY 4.0)}
\begin{document}
\maketitle

\noindent\emph{Source and licence.} A simple geometric proof for\\ the characterisation of e-merging functions. Version arXiv:2512.09708v3 (CC BY 4.0), distributed by arXiv under the Creative Commons Attribution 4.0 International licence, http://creativecommons.org/licenses/by/4.0/, as recorded in the arXiv licence metadata for this exact version and shown on the abstract page https://arxiv.org/abs/2512.09708v3. Full licence notice: \texttt{licenses/arxiv-2512.09708-CC-BY-4.0.txt}.\par\smallskip

\noindent\textit{Editorial scope.} Counted unit: the article's lemma lemma (source0.tex lines 84--86). The article's complete original proof of it is delivered with it (source0.tex lines 87--91, one \texttt{proof} environment of 5 lines). No mathematics is written, repaired or re-derived: the statement and the proof are the article's own lines, in the article's own notation, and no retained line is substituted or re-flowed.  No further source-local context is needed: every cross-reference of every retained line resolves inside the two delivered spans.  Nothing was omitted from any retained span, so the package carries no omission record.  No retained line contains a citation, so the wrapper carries no bibliography and no bibliography entry.  No retained line contains a figure environment, an image-inclusion command or drawing code, so no image is delivered and the package declares no asset dependency.  Editorial formatting only, in addition: an AMS \texttt{amsart} preamble that restates the article's own declarations and macros used by the retained lines, namely the environments \texttt{lemma} on separate counters, together with the standard packages those lines need.  The retained environments print in this document as Lemma 1 in order of appearance, whereas the article prints the same items with the numbering of its own full document.  The complete native source of the article as distributed in the arXiv e-print for this exact version is bundled unchanged as \texttt{sources/190754/source0.tex}.
\begin{lemma}\label{lemma}
    Fix an integer $n\geq 1$, consider $n$ vectors $\mathbf{u}^{(1)},\dots, \mathbf{u}^{(n)}\in [0,\infty)^K$, and a weight vector $(q_1,\dots, q_n)\in[0,1]^n$, such that $\sum_{j=1}^n q_j =1$. We have the following implication: $$\sum_{j=1}^n q_j \mathbf{u}^{(j)}= \mathbf{1}\qquad\implies\qquad\sum_{j=1}^n q_j F(\mathbf{u}^{(j)})\leq 1\,.$$
\end{lemma}

\begin{proof}
    Let $\X = \{x_1,\dots, x_n\}$ be a  set made of $n$ distinct elements, and consider the hypothesis $\Hh = \{Q\}$ on $\X$, with $Q = \sum_{j=1}^n q_j \delta_{x_j}$. Define $K$ random variables $E_1,\dots, E_K$ on $\X$ via $E_k(x_j) = u^{(j)}_k$. Then, 
    $$\E_Q[E_k] = \sum_{j=1}^n q_j E_k(x_j) = \sum_{j=1}^n q_j u^{(j)}_k = 1\,,$$ and so each $E_k$ is an e-variable for $\Hh$. We have
    $$\sum_{j=1}^n q_j F\big(\mathbf{u}^{(j)}\big) = \sum_{j=1}^n q_j F\big(E_1(x_j),\dots, E_K(x_j)\big) = \E_Q[F(E_1,\dots, E_K)] \leq1\,,$$ since $F$ being a generalised e-merging function implies that $F(E_1,\dots, E_K)$ is an e-variable for $\Hh$.
\end{proof}

\end{document}
